\section{Precision and resource demand as a finite effect calculus}

\subsection{Demand vectors}

A scalar precision is often insufficient.  A rational operation may need
numerator bits, denominator separation, and normalization fuel; a root routine
may need coefficient precision, disk radius, derivative gap, and branch depth.
We therefore use a fixed-dimensional demand vector
\[
      d=(d_0,\ldots,d_{k-1})\in\N^k
\]
ordered componentwise, with join given by componentwise maximum.  The dimension
\(k\) is fixed by the application, so tuple coding and all operations are
primitive recursive.

Each node has a monotone pullback map
\[
       T_{v\leftarrow u}:D_v\to D_u
\]
for every incoming edge \(u\to v\).  It states the demand on \(u\)'s output
needed to satisfy a demand on \(v\)'s output.

\subsection{Acyclic backward propagation}

Let the graph be topologically ordered.  Initialize every demand to zero and
set the terminal demand to the requested value.  Traverse nodes in reverse
order.  At node \(v\), join
\(T_{v\leftarrow u}(d_v)\) into each predecessor demand \(d_u\).

\begin{theorem}[Least acyclic demand]\label{thm:acyclic-demand}
For a finite acyclic graph with monotone edge pullbacks, the reverse pass
computes the least demand assignment satisfying the terminal request and every
edge constraint
\[
       d_u\ge T_{v\leftarrow u}(d_v).
\]
\end{theorem}

\begin{proof}
Soundness follows immediately when each edge is processed: the predecessor is
updated by a join containing the required value.  For minimality, let \(e\) be
any satisfying assignment.  Proceed in reverse topological order.  The terminal
demand computed by the algorithm is at most \(e\) by the terminal constraint.
If \(d_v\le e_v\), monotonicity gives
\(T_{v\leftarrow u}(d_v)\le T_{v\leftarrow u}(e_v)\le e_u\).
Every join inserted into \(d_u\) is therefore below \(e_u\), so the final
\(d_u\le e_u\).  This proves leastness.
\end{proof}

The algorithm is a single bounded reverse loop over a coded edge table.  Its
proof uses a quantifier-free table invariant and therefore fits
Theorem~\ref{thm:trace}.

\subsection{Why cyclic demands arise}

Cycles occur when two representations mutually constrain precision.  Examples
include:

\begin{itemize}
\item normalization of a rational expression needs a denominator gap, while
proving the gap needs a sharper approximation of the expression;
\item a chart transition requires coordinate precision determined by a local
Lipschitz bound, and the bound is selected from the chart reached at that
precision;
\item a Newton step needs derivative precision based on the current root disk,
while shrinking the disk uses the derivative estimate;
\item a network of mutually recursive finite observations shares error budgets.
\end{itemize}

The repair is not unbounded fixed-point iteration.  The theorem request already
imposes finite maximum resources.  Work in a finite product lattice and compute
the least closure there.

\subsection{Bounded least fixed points}

Let
\[
     L=\prod_{i<k}\{0,1,\ldots,B_i\}
\]
and let \(F:L\to L\) be monotone.  Define the inflationary closure operator
\(G(x)=x\join F(x)\) and the sequence
\[
       x_0=a,\qquad x_{t+1}=G(x_t).
\]
Here \(a\) is the externally requested minimum demand.

\begin{theorem}[Finite monotone closure]\label{thm:finite-closure}
The sequence stabilizes after at most
\[
      H=1+\sum_{i<k}(B_i-a_i)
\]
iterations.  Its stable value \(x^*\) is the least prefixed point of \(F\) above
\(a\):
\[
  a\le x^*,\quad F(x^*)\le x^*,
  \quad\text{and}\quad
  (a\le y\ \wedge\ F(y)\le y)\Rightarrow x^*\le y.
\]
\end{theorem}

\begin{proof}
The sequence is nondecreasing because \(x_t\le G(x_t)\).  Give the box the rank
\(\rank(x)=\sum_i x_i\).  If \(x_{t+1}\ne x_t\), some coordinate increases by
at least one, so \(\rank(x_{t+1})>\rank(x_t)\).  The rank can increase at most
\(\sum_i(B_i-a_i)\) times.  Thus a repeated value, hence a fixed point of
\(G\), occurs within the stated number of evaluated steps.  At a fixed point,
\(F(x^*)\le x^*\).

For leastness, let \(y\) be any prefixed point above \(a\).  Induct on \(t\).
The base \(x_0=a\le y\) holds.  If \(x_t\le y\), monotonicity gives
\(F(x_t)\le F(y)\le y\), and therefore
\(x_{t+1}=x_t\join F(x_t)\le y\).  Hence every iterate, including \(x^*\), is
below \(y\).
\end{proof}

Every quantity in this proof is a fixed-dimensional tuple operation.  The
iteration is primitive recursive with fuel \(H\), and the invariant
\[
 x_t\in L,\quad a\le x_t,\quad x_t\le x_{t+1},\quad
 \rank(x_t)\ge\rank(a)+\min(t,\text{number of strict steps})
\]
can be Booleanized.  This gives a direct path to an iPRA proof.

\subsection{Cycles with failure}

A demand transformer may detect that the finite box is too small: a required
precision exceeds \(B_i\), a denominator gap certificate is unavailable, or a
resource budget is exhausted.  Totalize the transformer with a failure flag.
The closure theorem then proves one of two finite outcomes:

\begin{enumerate}
\item a least closed demand inside the box; or
\item a certified failing node and coordinate, showing that the requested
contract cannot be discharged under the supplied resource envelope.
\end{enumerate}

This is preferable to silently clipping demands at the bound.  Clipping would
produce a total function but an invalid correctness theorem.

\subsection{Changing-parameter recursion}

Some recurrences change their own parameters.  For example, an approximation
at stage \(t+1\) may request stage-dependent precision \(p(t,s_t)\).  A naive
formalization attempts induction over a family whose formula changes with
\(t\), producing difficult substitutions.

The framework uses a backward-demand/forward-replay transformation.

\begin{construction}[Backward-demand/forward-replay]
For a fixed horizon \(T\):
\begin{enumerate}
\item run a reverse trace from the terminal request, computing the demand and
auxiliary parameters needed at every time;
\item code this finite schedule in one natural number;
\item run the original recurrence forward while reading the precomputed
schedule;
\item use one invariant saying that the schedule entry is valid and the
current state satisfies the corresponding contract.
\end{enumerate}
\end{construction}

The induction formula is now fixed: it refers to table lookup at the current
clock, not to a syntactically changing family.  The reverse pass is itself a
clocked PR trace and can be proved correct separately.

\begin{theorem}[Scheduled recurrence]
Assume each local step contract is monotone in its declared demand and is
proved in \(\KPrf\).  Then a finite changing-parameter recurrence with an
explicit horizon compiles to one reverse demand trace, one forward execution
trace, and a terminal \(\KPrf\) derivation.
\end{theorem}

\begin{proof}
The reverse trace produces demands satisfying every local edge constraint by
Theorem~\ref{thm:acyclic-demand} or
Theorem~\ref{thm:finite-closure} when local cycles are present.  During forward
replay, specialize each local step theorem at the scheduled demand and current
state.  The trace invariant supplies the step's assumptions.  At the host
proof-graph level, Theorem~\ref{thm:dag-compile} composes the reverse-trace
proof, local step proofs, and terminal finalizer.
\end{proof}

\subsection{Error budgets}

For approximate mathematics, demands are dual to error budgets.  If a node has
an output tolerance \(\varepsilon\), it computes input tolerances
\(\varepsilon_i\) and proves an inequality such as
\[
    E_{\mathrm{out}}
       \le c_0E_0+\cdots+c_rE_r+E_{\mathrm{round}}
       \le\varepsilon.
\]
All quantities are encoded rationals.  The inequality is decidable after
clearing positive denominators.  Polynomial parts may be discharged by the
ring backend; positivity and order side conditions remain explicit clauses.

A useful normalization is to store dyadic exponents instead of arbitrary
rational tolerances.  Addition of error terms then uses a conservative minimum
precision plus a small overhead.  The exact rational representation remains
available when sharp constants matter.

\subsection{Resource provenance}

Every bound should carry a provenance tag in the artifact metadata:

\begin{itemize}
\item \emph{structural}: derived from code size or a finite state count;
\item \emph{analytic}: derived from a modulus, Lipschitz constant, or gap;
\item \emph{algebraic}: derived from degree, coefficient height, or elimination
bounds;
\item \emph{external}: supplied as a theorem hypothesis or imported theorem;
\item \emph{experimental}: used only in tests and not in the proof statement.
\end{itemize}

The proof term need not mention the tag, but the research package should.  This
prevents an experimentally observed cutoff from being mistaken for a proved
height bound.
